\subsection{CHANGE OF RING OF SCALARS}

Let A, B be two rings and \(\rho\) a homomorphism of the ring B into the ring A. For every A-module E, the external law \((\beta, x) \mapsto \rho(\beta)x\) defines (with addition) a B-module structure said to be associated with \(\rho\) and the A-module structure on E; this B-module is denoted by \(\rho_{*}(\mathbf{E})\) or \(E_{[B]}\) (and even simply E if no confusion can arise). In particular, if B is a subring of A and \(\rho: B \to A\) is the canonical injection, \(E_{[B]}\) is called the B-module obtained by restricting the ring of scalars A to B; by an abuse of language, this expression is also used when the homomorphism \(\rho\) is arbitrary.

If \(\mathbf{F}\) is a submodule of the A-module \(\mathbf{E},\rho_{*}(\mathbf{F})\) is a submodule of \(\rho_{*}(\mathbf{E})\) and \(\rho_{*}(\mathbf{E} / \mathbf{F})\) is equal to \(\rho_{*}(\mathbf{E}) / \rho_{*}(\mathbf{F})\) .

Let E, F be two A-modules; every A-linear mapping \(u: E \rightarrow F\) is also a B-linear mapping \(E_{[B]} \rightarrow F_{[B]}\) denoted by \(\rho_{*}(u)\) ; in other words, there is a canonical injection of Z-modules

\[
\operatorname{Hom} _ {\mathbf {A}} (\mathbf {E}, \mathbf {F}) \rightarrow \operatorname{Hom} _ {\mathbf {B}} (\mathbf {E} _ {[ \mathbf {B} ]}, \mathbf {F} _ {[ \mathbf {B} ]}).\tag{47}
\]

This mapping is not necessarily bijective; in other words a B-linear mapping \(E_{[B]} \rightarrow F_{[B]}\) is not necessarily A-linear. For example, a sub-B-module of \(E_{[B]}\) is not necessarily a sub-A-module of E: if A is a field and B a subfield of A, the vector subspace \(B_{s}\) of the vector B-space \((A_{s})_{[B]}\) is not a vector sub-A-space if \(B \neq A\) .

It is immediate that, for every family \((\mathbf{E}_t)_{t\in I}\) of A-modules, the B-module \(\rho_{*}\left(\prod_{i\in I}\mathbf{E}_{i}\right)\left(\text{resp. } \rho_{*}\left(\bigoplus_{i\in I}\mathbf{E}_{i}\right)\right)\) is equal to \(\prod_{i\in I}\rho_{*}(\mathbf{E}_{i})\) (resp. \(\bigoplus_{i\in I}\rho_{*}(\mathbf{E}_{i})\) ).

Every generating system of \(\rho_{*}(\mathbf{E})\) is a generating system of E but the converse is not necessarily true.

PROPOSITION 24. Let A, B be two rings and \(\rho: \mathbf{B} \to \mathbf{A}\) a ring homomorphism.

\begin{enumerate}
\def\labelenumi{(\roman{enumi})}
\item
  If \(\rho\) is surjective, the canonical mapping (47) is bijective. For every A-module E, every sub-B-module of \(\rho_{*}(\mathbf{E})\) is a sub-A-module of E; every generating system of E is a generating system of \(\rho_{*}(\mathbf{E})\) .
\item
  If \(\rho\) is injective, every free family in the A-module \(\mathbf{E}\) is a free family in the B-module \(\rho_{*}(\mathbf{E})\) .
\end{enumerate}

The proposition follows immediately from the definitions.

Note that even if \(\rho\) is injective, a free family in \(\rho_{*}(\mathbf{E})\) is not necessarily free in \(\mathbf{E}\) .

*For example, 1 and \(\sqrt{2}\) do not form a free system in \(\mathbf{R}\) considered as a vector \(\mathbf{R}\) -space, although they form a free system in \(\mathbf{R}\) considered as a vector \(\mathbf{Q}\) -space (cf.~Remark 1).*

PROPOSITION 25. Let A, B be two rings, \(\rho: \mathbf{B} \to \mathbf{A}\) a ring homomorphism and E an A-module. Let \((\alpha_{\lambda})_{\lambda \in L}\) be a generating system (resp. free family of elements, basis) of A considered as a left B-module. Let \((a_{\mu})_{\mu \in M}\) be a generating system (resp. free family of elements, basis) of the A-module E. Then \((\alpha_{\lambda}a_{\mu})_{(\lambda,\mu) \in L \times M}\) is a generating system (resp. (when \(\rho\) is injective) free family of elements, basis) of the B-module \(\rho_{*}(\mathbf{E})\) .

If \(x = \sum_{\mu \in \mathbf{M}} \gamma_{\mu} a_{\mu}\) , where \(\gamma_{\mu} \in \mathbf{A}\) and \((\alpha_{\lambda})\) is a generating system of \(\mathbf{A}\) , we may write \(\gamma_{\mu} = \sum_{\lambda \in L} \rho(\beta_{\lambda \mu}) \alpha_{\lambda}\) , with \(\beta_{\lambda \mu} \in \mathbf{B}\) , for all \(\mu \in \mathbf{M}\) , whence \(x = \sum_{\lambda, \mu} \rho(\beta_{\lambda \mu}) \alpha_{\lambda} a_{\mu}\) . On the other hand, if \((\alpha_{\lambda})\) and \((a_{\mu})\) are free families, a relation \(\sum_{\lambda, \mu} \rho(\beta_{\lambda \mu}) \alpha_{\lambda} a_{\mu} = 0\) , with \(\beta_{\lambda \mu} \in \mathbf{B}\) , may be written \(\sum_{\mu \in M} \left( \sum_{\lambda \in L} \rho(\beta_{\lambda \mu}) \alpha_{\lambda} \right) a_{\mu} = 0\) ; it thus implies \(\sum_{\lambda \in L} \rho(\beta_{\lambda \mu}) \alpha_{\lambda} = 0\) for all \(\mu \in M\) and therefore \(\beta_{\lambda \mu} = 0\) for all \(\lambda, \mu\) if \(\rho\) is injective.

COROLLARY. If A is a finitely generated left B-module and E a finitely generated left A-module, \(\rho_{*}(\mathbf{E})\) is a finitely generated left A-module.

Let C be a third ring, \(\rho':\mathbf{C}\to\mathbf{B}\) a ring homomorphism and \(\rho''=\rho\circ\rho'\) the composite homomorphism. It follows immediately from the definitions that \(\rho_{*}^{\prime \prime}(\mathbf{E}) = \rho_{*}^{\prime}(\rho_{*}(\mathbf{E}))\) for every A-module E. In particular, if \(\rho\) is an isomorphism of B onto A, then \(\mathbf{E} = \rho_{*}^{\prime}(\rho_{*}(\mathbf{E}))\) , where \(\rho^{\prime}\) denotes the inverse isomorphism of \(\rho\) .

Remarks. (1) Let K be a field and A a subring of K with the following property: for every finite family \((\xi_{i})_{1\leqslant i\leqslant n}\) of elements of K, there exists a \(\gamma\in A\) which is non-zero and such that \(\gamma\xi_{i}\in A\) for \(1\leqslant i\leqslant n\) (a hypothesis which is always satisfied when A is commutative and K is the field of fractions of A). Let E be a vector space over K and \(E_{[A]}\) the A-module obtained by restricting the ring of scalars to A. Then, if a family \((x_{\lambda})_{\lambda\in L}\) is free in \(E_{[A]}\) it is also free in E. Attention may be confined to the case where \(L=\{1,n\}\) ; if there were a relation \(\sum_{i=1}^{n} \xi_i x_i = 0\) with \(\xi_i \in \mathbf{K}\) , the \(\xi_i\) not all zero, it would follow that for all \(\beta \in \mathbf{A}, \sum_{i=1}^{n} (\beta \xi_i) x_i = 0\) . By hypothesis we can suppose \(\beta \neq 0\) in \(\mathbf{A}\) such that \(\beta \xi_{i} = \alpha_{i}\) belongs to A for all \(i\) ; but the relation \(\sum_{i=1}^{n} \alpha_{i} x_{i} = 0\) is contrary to the hypothesis, the \(\alpha_{i}\) being not all zero.

\begin{enumerate}
\def\labelenumi{(\arabic{enumi})}
\setcounter{enumi}{1}
\tightlist
\item
  If the ring homomorphism \(\rho: \mathbf{B} \to \mathbf{A}\) is surjective and \(\mathfrak{b}\) is its kernel (so that \(\mathbf{A}\) is canonically identified with \(\mathbf{B} / \mathfrak{b}\) ), then, for every A-module \(\mathbf{E}\) , \(\mathfrak{b}\) is contained in the annihilator of \(\rho_*(\mathbf{E})\) and \(\mathbf{E}\) is the A-module derived from \(\rho_*(\mathbf{E})\) by the process defined in no. 12.
\end{enumerate}

Let A, B be two rings and \(\rho: \mathbf{B} \to \mathbf{A}\) a homomorphism. Let E be an A-module and F a B-module; a B-linear mapping \(u: \mathbf{F} \to \rho_{*}(\mathbf{E})\) (also called a B-linear mapping of F into E if no confusion arises) is also called a semi-linear mapping (relative to \(\rho\) ) of the B-module F into the A-module E; it is also said that the ordered pair \((u, \rho)\) is a dimorphism of F into E; this therefore means that, for \(x \in \mathbf{F}\) , \(y \in \mathbf{F}\) and \(\beta \in \mathbf{B}\) ,

\[
\left\{ \begin{array}{l} u (x + y) = u (x) + u (y) \\ u (\beta x) = \rho (\beta) u (x). \end{array} \right.\tag{48}
\]

The set \(\operatorname{Hom}_{\mathbf{B}}(\mathbf{F},\rho_{*}(\mathbf{E}))\) of B-linear mappings of \(\mathbf{F}\) into \(\mathbf{E}\) is also written as \(\operatorname{Hom}_{\mathbf{B}}(\mathbf{F},\mathbf{E})\) if no confusion can arise.

When \(\rho\) is an isomorphism of B onto A, the relation \(u(\beta x) = \rho(\beta) u(x)\) for all \(\beta \in B\) may also be written as \(u(\rho'(\alpha)x) = \alpha x\) for all \(\alpha \in A\) , where \(\rho'\) denotes the inverse isomorphism of \(\rho\) ; to say that \(u\) is semi-linear for \(\rho\) is then equivalent to saying that \(u\) is an A-linear mapping of \(\rho_*(F)\) into E.

Example. It has been seen (no. 1) that a homothety \(h_{\alpha} \colon x \mapsto \alpha x\) on an A-module \(\mathbf{E}\) is not necessarily a linear mapping. But if \(\alpha\) is invertible, \(h_{\alpha}\) is a semi-linear mapping (which is moreover bijective) relative to the inner automorphism \(\xi \mapsto \alpha \xi \alpha^{-1}\) of \(\mathbf{A}\) , for \(\alpha(\lambda x) = (\alpha \lambda \alpha^{-1})(\alpha x)\) .

Let C be a third ring, \(\rho':\mathrm{C}\to\mathrm{B}\) a homomorphism and G a C-module. If \(v:\mathrm{G}\to\mathrm{F}\) is a semi-linear mapping relative to \(\rho'\) , the composition \(w=u\circ v\) is a semi-linear mapping of G into E relative to the homomorphism \(\rho''=\rho\circ\rho'\) . If \(\rho\) is an isomorphism and \(u:\mathrm{F}\to\mathrm{E}\) is a bijective semi-linear mapping relative to \(\rho\) , the inverse mapping \(u':\mathrm{E}\to\mathrm{F}\) is a semi-linear mapping relative to the inverse isomorphism \(\rho':\mathrm{A}\to\mathrm{B}\) of \(\rho\) .

It is thus seen that, for the species of structure defined by giving on an ordered pair (A, E) of sets a ring structure on A and a left A-module structure on E, the dimorphisms \((u, \phi)\) can be taken as morphisms (Set Theory, IV, §2, no. 1); we shall always assume in what follows that this choice of morphisms has been made.

Remark (3). Let \(A_{1}\) , \(A_{2}\) be two rings, \(A = A_{1} \times A_{2}\) their product and let \(e_{1} = (1, 0)\) , \(e_{2} = (0, 1)\) in A, so that \(A_{1}\) and \(A_{2}\) are canonically identified with the two-sided ideals \(Ae_{1}\) and \(Ae_{2}\) of A. For every A-module E, \(e_{1}E\) and \(e_{2}E\) are sub-A-modules \(E_{1}\) , \(E_{2}\) of E, annihilated respectively by \(e_{2}\) and \(e_{1}\) , so that, canonically identifying \(A/Ae_{2}\) with \(A_{1}\) and \(A/Ae_{1}\) with \(A_{2}\) , \(E_{1}\) (resp. \(E_{2}\) ) is given an \(A_{1}\) -module (resp. \(A_{2}\) -module) structure. Moreover, E is the direct sum of \(E_{1}\) and \(E_{2}\) , for every \(x \in E\) can be written as \(x = e_{1}x + e_{2}x\) and the relation \(e_{1}x = e_{2}y\) implies \(e_{1}x = e_{1}^{2}x = e_{1}e_{2}y = 0\) . Conversely, for every ordered pair consisting of an \(A_{1}\) -module \(F_{1}\) and an \(A_{2}\) -module \(F_{2}\) , let \(E_{1}\) be the A-module \((p_{1}) * (F_{1})\) , \(E_{2}\) the A-module \((p_{2}) * (F_{2})\) , \(p_{1}\) and \(p_{2}\) being the projections of A onto \(A_{1}\) and \(A_{2}\) respectively; then in the A-module \(E = E_{1} \oplus E_{2}\) , \(E_{1} = e_{1}E\) , \(E_{2} = e_{2}E\) . The study of A-modules is thus reduced to that of \(A_{1}\) -modules and that of \(A_{2}\) -modules. In particular, every submodule M of E is of the form \(M_{1} \oplus M_{2}\) , where \(M_{1} = e_{1}M\) and \(M_{2} = e_{2}M\) .

